% Part: incompleteness
% Chapter: incompleteness-provability
% Section: introduction

\documentclass[../../../include/open-logic-section]{subfiles}

\begin{document}

\olfileid{inc}{inp}{int}

\olsection{مقدمة}

رأى هيلبرت أن نسقًا من البديهيات لبنية رياضية، مثل الأعداد الطبيعية،
لا يكون وافيًا ما لم يتيح اشتقاق جميع العبارات الصادقة عن البنية.
وحين يقترن هذا باهتمامه اللاحق بنظم الاستنباط الصورية، فإنه يوحي بأنه
كان يرى وجوب ضمان ألا تكون، مثلًا، النظم الصورية التي نستعملها
للاستدلال في الأعداد الطبيعية متسقة فحسب، بل \emph{تامة} أيضًا؛ أي إن
كل عبارة في لغتها تكون إما !!{derivable} أو يكون نفيها كذلك. وتبين
مبرهنة غودل الأولى لعدم الاكتمال أنه لا يوجد نسق بديهي من هذا القبيل:
فلا يوجد للحساب نسق صوري تام متسق و!!{axiomatizable}. بل لا تكون أي
نظرية رياضية متسقة و!!{axiomatizable} و«قوية بما يكفي» تامة.

وكان لهيلبرت هدف أهم، هو محور برنامجه لتسويغ الرياضيات الحديثة
(«الكلاسيكية»)، ويتمثل في إيجاد براهين تناهية على اتساق النظم الصورية
التي تمثل الاستدلال الكلاسيكي. ولذلك كانت مبرهنة غودل الثانية لعدم
الاكتمال ضربة أشد بكثير لبرنامج هيلبرت. ويمكن صياغة المبرهنة الثانية،
مثل الأولى، بعبارات تقريبية. فهي تقول، على وجه التقريب، إنه لا تستطيع
أي نظرية حساب متسقة وقوية بما يكفي إثبات اتساقها الخاص. وسيتعين أن نفهم
«قوية بما يكفي» على أنها تتضمن أكثر قليلًا من~$\Th{Q}$.

يمكن العثور على الفكرة الكامنة وراء برهان غودل الأصلي لمبرهنة عدم
الاكتمال في مفارقة إبمينيدس. فقد زعم إبمينيدس، وهو كريتي، أن جميع أهل
كريت كاذبون؛ والصورة الأصرح للمفارقة هي العبارة «هذه الجملة كاذبة».
وتمكن غودل، في جوهر الأمر، من إحلال !!{derivability} محل الصدق وصوغ
!!a{sentence} صوريًا تقرر، على نحو غير مباشر، أنها هي نفسها غير
!!{derivable}. فلو كانت تلك الـ!!{sentence} !!{derivable}، لصارت
النظرية غير متسقة. وبيّن غودل أن نفي تلك الـ!!{sentence} غير
!!{derivable} أيضًا من نسق البديهيات الذي كان ينظر فيه. (وفي هذا الجزء
الثاني كان على غودل أن يفترض أن النظرية~$\Th{T}$ «متسقة-$\omega$».
ويرتبط الاتساق-$\omega$ بالاتساق، لكنه خاصية أقوى.\footnote{أي إن كل
نظرية متسقة-$\omega$ متسقة، ولا يصح العكس.} وبعد غودل ببضع سنوات، بيّن
روسر أن افتراض مجرد اتساق~$\Th{T}$ يكفي.)

يتمثل التحدي الأول في فهم كيفية بناء !!a{sentence} تحيل إلى نفسها.
ولكل !!{formula}~$!A$ في لغة~$\Th{Q}$، ليرمز~$\gn{!A}$ إلى العدد
الاسمي المناظر لـ$\Gn{!A}$. تأمل معنى ذلك: $!A$ هي !!a{formula} في
لغة~$\Th{Q}$، و$\Gn{!A}$ عدد طبيعي، و$\gn{!A}$ \emph{حد} في
لغة~$\Th{Q}$. ومن ثم يكون لكل !!{formula}~$!A$ في لغة~$\Th{Q}$
\emph{اسم} هو~$\gn{!A}$، وهذا الاسم حد في لغة~$\Th{Q}$؛ وبذلك يتوافر
لنا إطار تصوري تستطيع فيه الـ!!{formula}s في لغة~$\Th{Q}$ أن «تقول»
أشياء عن !!{formula}s أخرى. وتعرف اللمة الآتية بلمّة النقطة الثابتة.

\begin{lem}
لتكن $\Th{T}$ أي نظرية تمتدّ~$\Th{Q}$، ولتكن $!B(x)$ أي !!{formula}
ليس فيها حرًا إلا المتغير~$x$. عندئذ توجد !!a{sentence}~$!A$ بحيث
$\Th{T}\Proves !A\liff !B(\gn{!A})$.
\end{lem}

تقرر اللمة أنه، متى أعطيت أي خاصية~$!B(x)$، وجدت !!a{sentence}~$!A$
تقرر أن «$!B(x)$ تصدق عليّ»، وأن $\Th{T}$ «تعلم» ذلك.

فكيف يمكننا بناء !!a{sentence} كهذه؟ تأمل الصيغة الآتية لمفارقة
إبمينيدس، وهي من وضع كواين:
\begin{quote}
«تنتج كذبًا إذا سبقتها صورتها المقتبسة» تنتج كذبًا إذا سبقتها صورتها
المقتبسة.
\end{quote}
لا تحيل هذه الجملة إلى نفسها مباشرة. إنها تقرر شيئًا عن الموضوعات
التركيبية الواقعة بين علامتي الاقتباس فحسب، وهي في ذلك على قدم المساواة
مع جمل مثل:
\begin{enumerate}
\item «روبرت» اسم لطيف.
\item «ركضتُ.» جملة قصيرة.
\item «تتألف من ثلاث كلمات» تتألف من ثلاث كلمات.
\end{enumerate}
لكن ماذا يحدث إذا أخذ المرء العبارة «تنتج كذبًا إذا سبقتها صورتها
المقتبسة»، ووضع قبلها صورة مقتبسة منها؟ نحصل عندئذ على الجملة
الأصلية!{} وخلاصة الأمر أن الجملة تقرر أنها كاذبة.

\end{document}
